% appendix_strategies.tex: the identification graveyard: every design applied to these data, what it identifies, and why the ones that do not identify the rate fall short.
\section{The Identification Graveyard: Designs That Do Not Identify the Rate}\label{appendix:strategies}

The paper's estimate of the rate is the within-borrower one of Section \ref{sec:rate_people}; the across-borrower term margin adjusted for the penalty and selection is the check on it (Appendix \ref{appendix:ladder}), the rule-based experiments of Section \ref{sec:quasi} are the check on the sign and the horizon of the far-payment response, and the decline of the payment margin over the life of the loan is compared with the stopping model's own (Appendix \ref{appendix:mcliquidity}). Table \ref{tab:strategies} records every identification design applied to these data, with the object each identifies and, where it does not identify the rate, the reason, so that the estimates can be compared with the alternatives and a reader who prefers one of them can see what it costs. Appendix \ref{appendix:strategies:three} states in full the three objects in the bureau data that look as if they should reveal a discount rate and do not, and Appendix \ref{appendix:strategies:designs} states in full the eight design-based comparisons the paper uses.
\subsection{The record}

\begin{footnotesize}
\begin{longtable}{@{}p{0.21\textwidth}p{0.18\textwidth}p{0.42\textwidth}p{0.13\textwidth}@{}}
\caption{Identification strategies for the discount rate: the record}\label{tab:strategies}\\
% replication: the record of the designs, written from the memos and results named in each row; no generator
\toprule
Strategy & Object it would identify & What happened & Status \\
\midrule
\endfirsthead
\toprule
Strategy & Object it would identify & What happened & Status \\
\midrule
\endhead
\bottomrule
\endfoot
\multicolumn{4}{@{}l}{\emph{Estimating the discount rate from the default margins}} \\
Two separate logits on $\log m$ and $\log L$, ratio $\gamma_L/(\gamma_L+\gamma_m)$ & the payment-versus-total sensitivity of default & $\log L = \log m + \log T$, so each coefficient absorbs the other margin; the ratio is held near one half by construction (0.46--0.58 across eight fixed-effect menus) and its stability was a property of the construction & not an instrument for the rate \\
Joint regression, closed form $x/(e^x-1)$ with the full term & $\rho$ from the ratio of the two elasticities under exponential discounting & the marginal-utility ratio, survival in the loan and the balance penalty enter; on the model's own simulations the closed form overstates the ratio by 2.5--3$\times$, and the lifetime outcome reflects exposure & not an instrument for the rate \\
Fixed-window outcomes with the joint regression & the same ratio free of exposure & the ratio changes sign across windows (negative at six months, positive over the life) and across score deciles; the term margin is two channels plus exposure, not one & kept as a finding (Section \ref{sec:termmargin}) \\
\midrule
\multicolumn{4}{@{}l}{\emph{Instruments and designs for the term margin}} \\
Macro-rate instruments (Treasury lanes) & exogenous variation in the payment at a given term & failed the first-stage and exclusion gates; the control-function payment coefficient is negative & not an instrument for the rate \\
Lender size-menu instruments & exogenous variation in the term & lender-years with larger loans serve safer borrowers; the exclusion fails and one failure enters both instrumented coefficients & not an instrument for the rate \\
Leave-out lender-menu instruments (examiner design) & the same & lenders are not randomly assigned; kept as a diagnostic of the fixed-effect design (Appendix \ref{appendix:ivvalidation}) & diagnostic \\
Menu regression discontinuity at the 72-month amount cutoff & the effect of a rule-driven longer term at a fixed total & first stage strong but small (log $T$ rises by 0.01); second stage imprecise (elasticity $-2$ to $-8$, s.e. 1.3--4); score and income borderline continuous & imprecise \\
Difference estimator $\text{FE} - [\text{FE}(S) - \text{IV}(S)]$ under equi-confounding & the fixed-effect estimate corrected by its gap to the instrumented one where both exist & the pre-test problem; the gap is transportable only if invariant across groups, which is the criterion the specification search then imposes & measured; the invariance criterion of Appendix \ref{appendix:bridge} \\
Within-borrower fixed effects (two or more loans) & the margins free of every stable trait & the closest-to-causal profile the bureau offers; it does not remove the lender's private information at origination, which is time-varying & the validation target \\
Within-borrower ratio of the term response to the payment response, mapped through the window formula & the rate, with selection into contract length removed by design and the scale cancelled, conditional on the penalty channel having run off by twenty-four months & on 13.5 million borrowers with two or more auto loans (65 million loans): $0.02$ pooled, $0.11$ on average across score deciles; no gradient in score or income; a gradient in age; state ratios outside the formula's range (Section \ref{sec:rate_people}) & the paper's estimate \\
The same ratio by origination year, the borrower effect re-estimated inside each year & the rate by period & inside one year the borrower effect is identified only from borrowers with two auto loans in that year, and the ratio is at or above the formula's ceiling in every year; the period design interacts the elasticities with the period and identifies the borrower effect across years & replaced by the period design of Section \ref{sec:rate_people} \\
Servicer-transfer linking across furnisher keys & the exit hazard net of transfers, the formula's $h$ & a trade that moves to a new servicer exits under one key and reappears under another; matched on consumer and opening date, transfers are 15 percent of auto exits and 18 percent in the first year, and the net hazard is the one the estimates use (Appendix \ref{appendix:data}) & measured \\
Selection on observables (Oster) & the direction of selection on the term & higher score, shorter term: observed selection runs against the negative early-default coefficient & supporting \\
\midrule
\multicolumn{4}{@{}l}{\emph{Separating the channels of the term margin}} \\
Deficiency and recourse contrasts (state rules) & the balance-scaled penalty channel & confirmed pooled on the ten-percent extract ($+0.36$, s.e. $0.13$, at twelve months on autos; $+3.9$, s.e. $1.5$, on mortgages); on the one-percent panel with covariates there are 303 no-pursuit defaults and no power by group & pooled only \\
Full model inversion by group (optimal stopping, income process free) & $(\rho, \sigma)$ per group from the window profile and default rate & the model's term elasticities are a third the size of the observed ones and $\rho$ goes to the grid edges & not an instrument for the rate \\
Proposition \ref{prop:sensT} ratio with beliefs fixed (Appendix \ref{appendix:formula}) & $\rho$ per group given $\lambda$, survival and the penalty share & every profile is fitted, but $\rho$ and the penalty share are collinear across windows; the confidence set spans the grid; recovery on simulations is monotone but biased & needs an anchor \\
Contract-length bins as the extra dimension & separation of $\rho$ from the penalty share through the term & the exit hazard, not $\rho$, governs how the marginal payment's weight falls with the term; no separation & not an instrument for the rate \\
Charge-off amounts as the measured penalty base & the penalty channel's level and its shape over the loan's life, by group & measured with precision (elasticity to the payment 1.04; to the term 0.33 at six months rising to 1.79 after three years); under the first-order formula the measured shape makes both channels fall with the window while the observed ratio rises; the full model of Appendix \ref{appendix:mcliquidity} produces a rising ratio but at a third of the observed magnitude & the penalty shape is measured; the profile's rise is not explained \\
The equity margin on first mortgages: default against the house-price change since origination, against the payment (Section \ref{sec:mainresults}) & the ratio of a lump to a twenty-five-year stream, and its profile over loan age & 1.86 million mortgages, 191,581 events: elasticity $-2.5$, stronger without recourse, falling from $-5.9$ to $-1.1$ with age, convex in negative equity; in the model with a house, with i.i.d.\ income and a payment elasticity three times the measured one, the ratio to the payment response is between $-0.92$ and $-1.09$ at every rate, because the margin is the sale option (conditional on that calibration, Appendix \ref{appendix:persistent}) & measured; consistent with the model; not a rate \\
The term margin corrected rung by rung (Appendix \ref{appendix:ladder}) & $\rho$ from the ratio of the term to the payment response within twelve months, after the penalty and selection corrections & pooled charge-off: raw ratio 0.01 (rate near one), penalty correction to 0.81, selection correction to 0.39 (0.31), mapped to 0.04 with a set that spans the grid; ordering by group at the boundaries & measured; the check on the within-borrower estimate \\
\midrule
\multicolumn{4}{@{}l}{\emph{Timing variation inside the bureau}} \\
Deferred first payments and balloons & a rescheduling at a fixed balance & a few hundred auto loans; personal-installment deferrals are student-loan-like & unavailable \\
Pandemic forbearance & a rescheduling at scale & visible on mortgages only (1.3 percent at the 2020 peak), borrower-elected, and reported as current while in effect & not rule-based \\
Movers (place versus person) & the share of default variation due to place & answers its own question (slope 0.075, s.e. 0.013); it is not a discount-rate instrument & supporting \\
\midrule
\multicolumn{4}{@{}l}{\emph{What identifies $\rho$}} \\
Rule-based experiments, mapped through Proposition \ref{prop:sensT} (Section \ref{sec:quasi}) & $\rho$ for the populations the experiments treat & rescheduling against write-down fits best at zero and, aggregated over its five-year window, accepts $[0, 1.20]$ at five percent; the ARM resets anchor the near-horizon response; the forgiveness-versus-reduction pair is uninformative alone & measured; the check on the sign and the horizon \\
The payment margin over the life of the loan on the estimation sample & the run-off of the remaining payments' weight, by group & the run-off was large (the elasticity fell by two thirds over the contract); the sample's last-payment field, which reached the merged file through a separate charge-off extraction, dates the same cohorts' charge-offs three times as often in the first six months as the monthly panel does and a third as often late, so the profile was an artifact of default dating & not used; the monthly panel dates defaults from ratings \\
The payment margin over the life of the loan on the monthly panel, mapped through the first-order formula & the same, with survival and the marginal-utility ratio held constant & measured with standard errors of 0.02 over seven loan-age bins on 7.6 million loans; on the structural model's own simulations the formula attributes a run-off that the default option produces at any $\rho$ to a high $\rho$ (Appendix \ref{appendix:mcliquidity}), so the formula is not the mapping & measured; mapped through the model \\
The same profile, compared with the model's own run-off (Appendix \ref{appendix:mcliquidity}) & the set of rates at which the model reproduces each group's profile shape & the set spans the grid in nearly every group: the run-off validates the stopping problem and does not pin the rate & measured; not an instrument for the rate, by the option \\
\midrule
\multicolumn{4}{@{}l}{\emph{Timing variation inside the bureau}} \\
Adjustable-rate resets, 2004--08 originations, found from the payment path (Appendix \ref{sec:armreset}) & $D(\tau)\lambda$ at one to twelve months from the anticipation of a reset & 88,151 payment changes at hybrid reset ages, of which 9,224 are rate changes by the amortization identity and the rest escrow re-analyses, interest-only loans, recasts and closing trades; the index moved by less than 0.2 points across the 2010--13 reset years (first-stage $F$ below two); delinquency rises before a larger coming cut in every sample, a sign no discount function produces; 1,066 events on the clean cuts & the classification method and the pre-reset sign; no rate \\
The reset-type contrast: five-year against seven-year hybrids of the same vintage at the five-year reset month & the same, with vintage-specific trends removed & the pre-reset sign is the same within type, and the clean sample is too small to split & not built \\
The October 2023 student-loan payment resumption as a cross-loan experiment (Appendix \ref{sec:studentloan}) & the anticipation response at four to one months and the at-date response of auto delinquency to a known bill on another account, by its size & 438,311 holders, 743,487 auto loans, 17,711 events; the coefficient per \$100 of bill is negative in every month from March 2023, before the statute, through 2025 (composition of the exposure); the change at the bill is $-2.8$ ($0.6$) against $+5.9$ for an enforced payment of that size; half the bills unpaid and none reported delinquent for a year & measured; not an enforced payment, so not the design's premise \\
The same design on the enforced payments after May 2025 & the same & the coefficients from May to December 2025 do not separate from the preceding drift & open as the panel lengthens \\
Adjustable-rate resets in the ten-percent archive & the same as the one-percent resets with ten times the loans & the archive is observed twice a year through 2017, which does not resolve a monthly anticipation profile & not built, by the data's frequency \\
HAMP step-ups, 2014--19 & the same, from rule-fixed annual payment increases on modified loans & the amortization identity finds 2,478 rate rises of one half to one and a half points on loans below 4.5 percent in 2013--18, 69 of them twelve months after a previous one: fewer events than the reset cuts & too few events \\
Balloon payments, all families (\texttt{balloonpaymentamount}, \texttt{dateballoonpaymentdue}) & $D(\tau)\lambda$ at one to five years from the approach of a contractual lump sum & about 2,000 open auto and 3,000 dated mortgage balloons per one-percent snapshot, a few hundred delinquency events over their lives & too few loans \\
Credit-card promotional expirations & $D(\tau)\lambda$ at twelve to eighteen months for revolvers at the minimum & revolving trades' monthly minimum payments & not built \\
Contract-length bins as step responses within borrower (the discount kernel from the term margin) & the weight on payments at each added lag, from the step in default across contract lengths & steps are not monotone in the lag: the standard terms of 60 and 72 months stand above 48, 66 and 84 at every horizon, so the bins measure which product and lender the borrower took, not the weight on the added payments; the one clean pair, 60 against 72 months, is one point, not a curve & menu effects; not a kernel \\
Credit-card promotional-rate expirations & $D(\tau)\lambda$ at twelve to eighteen months from a payment increase fixed by the card agreement & delinquency rises by $0.033$ percentage points per one percent of the coming payment increase in the months before expiry and does not move at expiry; without cards of the same age without a promotion the anticipation cannot be separated from card age & measured; no control group \\
Pandemic forbearance exits & the rate at long horizons from deferred balances against repayment plans & narrative codes and payment paths & not built \\
\end{longtable}
\end{footnotesize}

Three lessons organize the table. First, a ratio of two coefficients on the same outcome cancels only symmetric confounding; every mechanism found loads on one margin, so sharing the outcome removed nothing. Second, the term margin across borrowers is the wrong instrument for a discount rate in these data: it bundles a preference channel bounded above by one, a balance-scaled penalty whose shape is measured, and a component that rises with the window and with the score, which neither channel produces at the observed magnitude; screening at origination on the lender's private information, whose predictive power decays, is the candidate the within-borrower design cannot remove. Third, what identifies a discount rate without a model of the penalty is variation in the timing of payments at a fixed balance, which the bureau does not contain on auto loans and the rule-based experiments do; the bureau's within-borrower ratio gives a rate with selection removed by design and the penalty channel assumed to have run off by twenty-four months, and the bureau's other contribution is the population-scale measurement of the objects the formula uses, and their heterogeneity across people and time.

\subsection{The three objects that do not identify the rate, in full}\label{appendix:strategies:three}

Three objects in the credit-bureau data look as if they should reveal a discount rate and do not, each for a reason that is itself a result of the theory. This subsection states each in full.

\paragraph{The response to a payment due now against payments due later cannot separate impatience from the present-payment premium.} Every response of default to a payment at date $s$ has the weight $D(s-1)Q_s\lambda$, and under quasi-hyperbolic discounting the present-bias parameter and the marginal-utility ratio appear in that weight only as the product $\beta\lambda$ (Theorem \ref{thm:betalambda}). This is not a limitation of a particular design: any experiment that moves a payment due now against payments due later, whether a payment holiday, a step-up, a rescheduling or the run-off of a contract, produces ratios that depend on $(\beta, \lambda)$ through their product, so the same data are consistent with a fully present-biased borrower whose marginal utility is the same in every state and with a time-consistent borrower who is liquidity-constrained at the margin. What separates them is a comparison of two dollars at the same date valued at different marginal utilities, a liquid dollar and an illiquid one, which the cash-on-hand and seizable-equity responses of \citet{indarte_moral_2023} provide (Corollary \ref{cor:separate}). The paper therefore measures the premium from that pair, uses the profile of equation \eqref{eq:lambda_p}, through Indarte's $0.2$, for the experiments' populations and for the bureau's groups, and reports the exponential rate beyond the first payment.

\paragraph{The response to the contract's length across borrowers cannot identify the rate on its own.} A longer contract at a fixed payment adds far payments, which a patient borrower weighs more heavily, but it also raises the balance the lender can pursue after default and the number of months in which a default can occur. Proposition \ref{prop:termmargin} writes the term response as the sum of the three: a preference channel bounded above by one in units of the payment response, a balance-scaled penalty channel that is negative and discounted at the contract rate, and, for lifetime outcomes, exposure; across borrowers it adds selection into contract length. Two features of the data show that the penalty and exposure dominate at short horizons. Within lender, place, year and score-by-income cell, a longer contract at a fixed payment lowers early delinquency for low-score borrowers and raises it for high-score borrowers, and the ratio of the term response to the payment response exceeds one above the fifth score decile, outside the preference channel at any rate (Section \ref{sec:termmargin}); and the early effect is weaker, by a third at twelve months on auto loans and by roughly half on mortgages, in the states whose laws bar the lender from pursuing the balance after repossession. The across-borrower term margin is reported as a finding about the cost of default, and it is mapped to a rate only after the penalty and selection are subtracted (Appendix \ref{appendix:ladder}); within borrower, with selection removed by design, it is taken at twenty-four months conditional on the penalty channel having run off (Section \ref{sec:rate_people}).

\paragraph{The run-off of the payment margin over the life of the loan is not mapped to a rate.} As a loan ages the remaining payment stream shortens, and a first-order formula, with survival and the premium held constant, would have the response to the payment decline at a rate governed by $\rho$. The optimal-stopping model says otherwise. A borrower who continues buys one period of standing plus the option to reconsider (Proposition \ref{prop:interiorT}), so the far payments enter her decision only with the weight of the states in which she will still be paying them, and when default is a threshold crossing in income that weight is small at every rate. On the model's own simulations, calibrated to the sample's default rate with exits at the measured hazard, the run-off ratio is between $0.77$ and $0.82$ at every rate from zero to $1.2$ (Appendix \ref{appendix:mcliquidity}), and the first-order formula applied to those simulations returns rates near $3$ at true rates of $0.03$ to $0.10$. The measured run-off, $0.75$ under the charge-off definition and $0.67$ under the ninety-day definition, lies below that range as a point estimate, and the 95 percent bands of the observed profiles overlap it in nearly every group (Section \ref{sec:mainresults}). The result is conditional on the calibration. The model reaches the sample's default rate only with a payment elasticity near five against the data's $0.45$ and with $0.93$ of lifetime default in the first year against the data's $0.13$; a version with a persistent income process, calibrated to the lifetime rate and to the first-year share, matches the timing of default, keeps the same elasticity at every persistence and dispersion, and has a run-off that is again insensitive to the rate but at a level of $0.2$ with a profile that declines from the first year, against the data's flat profile (Appendix \ref{appendix:persistent}). Neither version reaches the level of the payment margin, so the run-off is not used for the rate in either direction: it is not mapped through the first-order formula, and the model's insensitivity result is not a demonstration that the data's run-off contains no information about the rate.

\subsection{The design-based comparisons in full}\label{appendix:strategies:designs}

An estimate is design-based when the variation it uses is fixed by something the borrower does not choose, a state law, a statute, the date at which a neighbor bought a house, or when the borrower herself is held fixed, so that the comparison is between two of her own loans; an estimate is fixed effects plus adjustment when it compares loans that are alike on what the data record and is causal under the assumption of Section \ref{sec:identification}. The eight design-based comparisons are the following.

\begin{enumerate}
\item \emph{The payment margin within borrower.} The design compares a borrower's default on one of her auto loans with her default on another, within origination-year cells, on the ten-percent archive. Charge-off within twelve months responds to the monthly payment with an elasticity of $0.44$ (standard error $0.005$), and within twenty-four months $0.53$ ($0.003$) (Table \ref{tab:main_margins}, lower panel). This is the response of default to the payment due now for a fixed person. What the design leaves is what changed about her between the two loans that the lender saw and the cell does not record; Table \ref{tab:strategies} names that as the one confounder the within-borrower design cannot remove.

\item \emph{The term margin within borrower, by score.} The same design, with the contract's length at a fixed payment as the treatment, run decile by decile on the multi-loan borrowers with a score, the decile fixed by the borrower's first score, within borrower and origination-year cells (Table \ref{tab:within_by_group}). Within twenty-four months the elasticity is $0.30$ ($0.01$) pooled and between $0.19$ and $0.43$ across deciles, and its ratio to the within-borrower payment elasticity, the paper's sufficient statistic, is between $0.32$ and $0.53$ in every decile with standard errors from $0.02$ to $0.10$ (Table \ref{tab:rho_within}); within twelve months the elasticity is $-0.31$ ($0.01$). Across borrowers, within the same lender-year cells, the same elasticity runs from $-0.39$ in the bottom decile to $1.21$ in the top, so the whole gradient is selection into contract length, positive where lenders screen with maturity and negative where prime borrowers choose long terms for themselves \citep{hertzberg_screening_2018}. On personal installment loans the same design gives $0.28$ ($0.01$) at twenty-four months, and $0.30$, $0.29$ and $0.19$ by score tercile. The response of default to payments added at the end of a contract, net of every stable trait, is the same for every score in two loan families, and the selection correction of Appendix \ref{appendix:ladder} is measured, with a standard error of $0.01$ pooled, rather than assumed.

\item \emph{The deficiency-judgment contrast.} In sixteen states the law bars the lender from pursuing the unpaid balance after repossession on a loan below a threshold. The design compares auto loans below the threshold in those states with loans of the same size in states without a bar, within zip3-by-year cells (Table \ref{tab:term_deficiency}). Where the balance cannot be pursued, the term elasticity of default within twelve months is higher by $0.36$ ($0.13$), and within thirty-six months by $1.02$ ($0.27$); on mortgages, using recourse and non-recourse states, the difference is $3.9$ ($1.5$) at twelve months. This is the penalty channel of Proposition \ref{prop:termmargin} measured at its deficiency part: the cost of default rises with what is owed, by law, and a longer contract at the same payment lowers early default for that reason.

\item \emph{The equity margin.} Two households in the same zip code in the same quarter who bought their houses in different years hold different equity because prices moved between the two purchases; with zip-by-quarter and origination-year effects, the comparison is between origination cohorts under the same local conditions (Table \ref{tab:equity_margin}). The quarterly hazard of first 90-day delinquency responds to the log house-price change since origination with a coefficient of $-2.23$ ($0.09$), times one hundred, an elasticity of $-2.54$; the ratio of that response to the payment response is $-1.71$ ($0.07$), and it is $-2.38$ ($0.16$) in the states without deficiency judgments against $-1.54$ ($0.09$) with them. The response falls with the loan's age, from an elasticity of $-5.90$ in the first three years to $-1.07$ after seven, and the hazard is convex in negative equity, adding $3.9$ percentage points per quarter, four and a half times the base, once the price has fallen by half. These are the signatures of a walk-away decision and its dependence on recourse. The model with a house, with i.i.d.\ income and a payment elasticity of $4.7$ to $5.5$ against the measured $1.5$, returns a ratio between $-0.92$ and $-1.09$ at every rate (Table \ref{tab:equity_model}; the calibration's limits are those of Appendix \ref{appendix:persistent}), so the margin is consistent with the stopping problem and is not used for the rate.

\item \emph{The adjustable-rate resets.} The amortization identity finds, without a contract flag, $6{,}348$ rate cuts at hybrid reset ages among $429{,}078$ payment changes on 2004--08 first mortgages, with loan and vintage-by-calendar-month effects (Appendix \ref{appendix:timing}), and $11{,}859$ cuts when only the amortization before the cut is required. On that larger sample the response of first 90-day delinquency at the cut is $-0.50$ ($0.36$) per unit of log payment, times one hundred, and in the four months before it $-0.47$ ($0.13$): delinquency is higher before a larger coming cut. The sign is one that no discount function produces, and with the index flat across the reset years and $1{,}066$ events, no rate is estimated from this population (Section \ref{sec:mainresults}).

\item \emph{The student-loan resumption.} A statute fixed the date on which student-loan payments resumed, and the size of each bill was fixed by the student loan; the design follows auto default on the same consumer (Appendix \ref{appendix:timing}). The coefficient of first 90-day auto delinquency on the resumed bill, per hundred dollars, changes by $-2.8$ ($0.6$), in units of $10^{-5}$ per month, between the months before the statute and the months of the first bills, against $+5.9$ for an enforced payment of that size; the response is measured to within $\pm 1.4 \times 10^{-5}$, two percent of the base hazard, and is not positive at any horizon in any group. Auto default does not respond to a bill on another account that half of its recipients did not pay and that had no consequence for a year (Section \ref{sec:mainresults}).

\item \emph{The menu regression discontinuity.} Lenders admit a 72-month contract only above a loan-amount cutoff, and the design compares loans just above with loans just below it, within lender-year, at bandwidths of two to five thousand dollars. The log contract length rises by $0.011$ ($0.001$) at the cutoff, a first stage with an $F$ statistic near $200$ but small; the implied elasticity of default within twelve months to the term at a fixed total is $-2.6$ ($2.0$) at the middle bandwidth and $-2.1$ ($1.3$) at the widest. A longer term at a fixed total is a lower payment, and default falls, which is the sign of the payment margin; the estimate is not distinguishable from zero and is not used.

\item \emph{Movers.} Borrowers whose three-digit zip code changes between two of their loans identify how much of the geographic gradient in default is the place: the change in a mover's own default, regressed on the change in the stayers' default rate between destination and origin, has a slope of $0.075$ ($0.013$) over the life of the loan (Table \ref{tab:person_place}). The geography of Section \ref{sec:heterogeneity} is composition to that extent, and this is not a discount-rate object.
\end{enumerate}
